\scanpage{179}
\mspage{128}
($a)$ et $c)$ n'ayant d'ailleurs pas changé et $b$ devenant $b')$ $f'_1:\mathscr X_{/1}\to\mathscr X'_{/1}$ et $f'_2:\mathscr X_{/2}\to\mathscr X'_{/2}$ sont dans $W$). Alors la conclusion \inserted{du cor.} est équivalente à $W(7)$ \inserted{(formulation \unsure{duale|dure})} dans le cas du nouveau diagramme cubique, \unsure{qed.|q.e.d.}

3.2. \uline{Comparaison de}
\diagram{\begin{tikzcd}[ampersand replacement=\&,row sep=tiny,column sep=small]
\&X_1\\
\int X_0\arrow[ru]\arrow[rd]\&\\
\&X_2
\end{tikzcd}}
\uline{avec une somme amalgamée ordinaire}. \uline{Carrés $W$-homotopiquement cocartésiens}.

La catégorie intégrale
\diagram{\begin{tikzcd}[ampersand replacement=\&,row sep=tiny,column sep=small]
\&X_1\\
\mathscr X=\int X_0\arrow[ru]\arrow[rd]\&\\
\&X_2
\end{tikzcd}}
est solution du problème universel suivant, \deleted{\unsure{2}} variante 2-catégorique de celui que définit la limite inductive du diagramme : il s'agit, pour $Y$ objet variable dans $\mathit{Cat}$, de considérer les systèmes de trois foncteurs $f_0,f_1,f_2$ et de deux hom. fonctoriels $\lambda_1,\lambda_2$ :
\diagram{\begin{tikzcd}[ampersand replacement=\&,column sep=small,row sep=small]
\&\&X_0\arrow[dl,"\alpha_1"']\arrow[dr,"\alpha_2"]\arrow[dd,"f_0", ""{name=C,near end}]\&\\
(\ast)\&X_1\arrow[dr,"f_1"',""{name=L,near end}]\&\&X_2\arrow[dl,"f_2", ""{name=R,near end}]\\
\&\&Y\&\\
\arrow[from=L,to=C,"\lambda_1",shorten <=1pt,shorten >=2pt]
\arrow[from=R,to=C,"\lambda_2"',shorten <=1pt,shorten >=2pt]
\end{tikzcd}}
\[
\begin{aligned}
\lambda_1 &: f_1\alpha_1\to f_0\\
\lambda_2 &: f_2\alpha_2\to f_0
\end{aligned}
\]
% archive sheet number: 179
